\chapter[Qualità ripugnanti dei requisiti]{Riflessione sulle qualità repulsive dei requisiti}

\firstline{Yathā paccayaṁ pavattamānaṁ}

\begin{leader}
  [Ha꜓nda mayaṁ dhātu-paṭikūla-paccavekkhaṇa-pāṭhaṁ bhaṇāmase]
\end{leader}

\begin{english}
  [Cantiamo ora la riflessione sulle qualità ripugnanti dei requisiti]
\end{english}

%\suttaref{Trad.}%
[Yathā꜓ pa꜕ccayaṁ] pava꜓tt꜕amānaṁ dhātu꜕-ma꜓tta꜕m-ev'etaṁ

\begin{english}
Composti da soli elementi secondo cause e condizioni
\end{english}

Yad i꜓daṁ cī꜓varaṁ ta꜕d upa꜕bhuñja꜓ko c꜕a pu꜕gga꜕lo

\begin{english}
Sono queste vesti e così è la persona che le indossa;
\end{english}

Dhātu-ma꜓tta꜕ko

\begin{english}
Semplici elementi,
\end{english}

Ni꜓ssa꜕tto

\begin{english}
Non un essere,
\end{english}

Ni꜓jjīvo

\begin{english}
Senza vita autonoma,
\end{english}

Su꜓ñño

\begin{english}
E vuoti di un sé.
\end{english}

S꜕abbāni pa꜕na imāni cī꜓varāni a꜕jigu꜓ccha꜕nīyāni

\begin{english}
Nessuna di queste vesti è per sua natura disgustosa
\end{english}

\clearpage

Imaṁ pūti꜓-kāyaṁ pa꜕tvā

\begin{english}
Ma toccando questo corpo impuro
\end{english}

A꜕tiviya jigu꜓ccha꜕nīyāni jāyanti

\begin{english}
Diventano davvero ripugnanti.
\end{english}

Yathā꜓ pa꜕ccayaṁ pava꜓tt꜕amānaṁ dhātu꜕-ma꜓tta꜕m-ev'etaṁ

\begin{english}
Composto di soli elementi secondo cause e condizioni
\end{english}

Yad i꜓daṁ piṇḍa꜓pāto ta꜕d upa꜕bhuñja꜓ko c꜕a pu꜕gga꜕lo

\begin{english}
È questo cibo offerto e così la persona che lo consuma;
\end{english}

Dhātu-ma꜓tta꜕ko

\begin{english}
Semplici elementi,
\end{english}

Ni꜓ssa꜕tto

\begin{english}
Non un essere,
\end{english}

Ni꜓jjīvo

\begin{english}
Senza vita autonoma,
\end{english}

Su꜓ñño

\begin{english}
E vuoti di un sé.
\end{english}

S꜕abbo pa꜕nāyaṁ piṇḍa꜓pāto a꜕jigu꜓ccha꜕nīyo

\begin{english}
Niente di questo cibo è per sua natura disgustosa
\end{english}

\clearpage

Imaṁ pūti꜓-kāyaṁ pa꜕tvā

\begin{english}
Ma toccando questo corpo impuro
\end{english}

A꜕tiviya jigu꜓ccha꜕nīyo jāyati

\begin{english}
Diventa davvero ripugnante.
\end{english}

Yathā꜓ pa꜕ccayaṁ pava꜓tt꜕amānaṁ dhātu꜕-ma꜓tta꜕m-ev'etaṁ

\begin{english}
Composta da soli elementi secondo cause e condizioni
\end{english}

Yad i꜓daṁ senā꜓sanaṁ ta꜕d upa꜕bhuñja꜓ko c꜕a pu꜕gga꜕lo

\begin{english}
È questa dimora e così la persona che la abita;
\end{english}

Dhātu-ma꜓tta꜕ko

\begin{english}
Semplici elementi,
\end{english}

Ni꜓ssa꜕tto

\begin{english}
Non un essere,
\end{english}

Ni꜓jjīvo

\begin{english}
Senza vita autonoma,
\end{english}

Su꜓ñño

\begin{english}
E vuoti di un sé.
\end{english}

S꜕abbāni pa꜕na imāni senā꜓sanāni a꜕jigu꜓ccha꜕nīyāni

\begin{english}
Niente di questa dimora è per sua natura disgustosa
\end{english}

\clearpage

Imaṁ pūti꜓-kāyaṁ pa꜕tvā

\begin{english}
Ma toccando questo corpo impuro
\end{english}

A꜕tiviya jigu꜓ccha꜕nīyāni jāyanti

\begin{english}
Diventa davvero ripugnante.
\end{english}

Yathā꜓ pa꜕ccayaṁ pavatt꜕amānaṁ dhātu꜕-matta꜕m-ev'etaṁ

\begin{english}
Composti da soli elementi secondo cause e condizioni
\end{english}

Yad i꜓daṁ gi꜕lāna-pacca꜕ya꜕-bhesajja-pa꜕rikkhāro ta꜕d upa꜕bhuñja꜓ko c꜕a pu꜕gga꜕lo

\begin{english}
Sono questi rimedi curativi e così è la persona che li consuma;
\end{english}

Dhātu-ma꜓tta꜕ko

\begin{english}
Semplici elementi,
\end{english}

Ni꜓ssa꜕tto

\begin{english}
Non un essere,
\end{english}

Ni꜓jjīvo

\begin{english}
Senza vita autonoma,
\end{english}

Su꜓ñño

\begin{english}
E vuoti di un sé.
\end{english}

S꜕abbo pa꜕nāyaṁ gi꜕lāna-pacca꜕ya꜕-bhesajja-pa꜕rikkhāro a꜕jigu꜓ccha꜕nīyo

\begin{english}
Niente di questi rimedi curativi è per sua natura disgustosa
\end{english}

Imaṁ pūti꜓-kāyaṁ pa꜕tvā

\begin{english}
Ma toccando questo corpo impuro
\end{english}

A꜕tiviya jigu꜓ccha꜕nīyo jāyati

\begin{english}
Diventano davvero disgustosi.
\end{english}

\chapter{Riflessione sull'impermanenza}

\firstline{Sabbe saṅkhārā aniccā}

\begin{leader}
  [Handa mayaṁ aniccānussati-pāṭhaṁ bhaṇāmase]
\end{leader}

\begin{english}
  [Cantiamo ora la riflessione sull'impermanenza]
\end{english}

[Sa꜕bbe sa꜓ṅkhā꜓rā a꜕ni꜓ccā]

\begin{english}
Tutte le cose condizionate sono impermanenti;
\end{english}

Sa꜕bbe sa꜓ṅkhā꜓rā du꜕kkhā

\begin{english}
Tutte le cose condizionate sono dukkha;
\end{english}

Sa꜕bbe dhammā a꜕na꜓ttā

\begin{english}
Ogni cosa è vuota di un sé.
\end{english}

A꜕ddhuvaṁ jīvi꜓taṁ

\begin{english}
La vita non è sicura;
\end{english}

Dhuvaṁ ma꜓ra꜕ṇaṁ

\begin{english}
La morte è sicura;
\end{english}

A꜕vassaṁ mayā mari꜓ta꜕bbaṁ

\begin{english}
È inevitabile che io muoia;
\end{english}

Ma꜕raṇa-pa꜕riyosā꜓naṁ me jīvi꜓taṁ

\begin{english}
La morte è il culmine della mia vita;
\end{english}

Jīvitaṁ me ani꜓ya꜕taṁ

\begin{english}
La mia vita è incerta;
\end{english}

Maraṇaṁ me ni꜓ya꜕taṁ

\begin{english}
La mia morte è certa.
\end{english}

Vata

\begin{english}
Davvero,
\end{english}

A꜕yaṁ kāyo

\begin{english}
Questo corpo
\end{english}

A꜕ciraṁ

\begin{english}
Presto
\end{english}

A꜕peta꜕-viññāṇo

\begin{english}
Sarà privo di coscienza
\end{english}

Chu꜕ddho

\begin{english}
E gettato via.
\end{english}

A꜕dhise꜓ssa꜕ti

\begin{english}
Giacerà
\end{english}

Pa꜕ṭha꜕viṁ

\begin{english}
Sulla terra
\end{english}

Ka꜕liṅga꜓raṁ i꜕va

\begin{english}
Proprio come un tronco marcio,
\end{english}

Ni꜕ratthaṁ

\begin{english}
Completamente privo di utilità.
\end{english}

Aniccā vata sa꜓ṅkhā꜓rā

\begin{english}
In verità le cose condizionate non possono durare,
\end{english}

U꜕ppāda-vaya-dha꜓mmi꜕no

\begin{english}
La loro natura è sorgere e cadere,
\end{english}

U꜕ppajjitvā nirujjh꜓anti

\begin{english}
Essendo sorte, le cose devono cessare,
\end{english}

Tesa꜓ṁ vūpa꜕sa꜕mo sukho

\begin{english}
La loro quiete è vera felicità.
\end{english}

\chapter{Veri e falsi rifugi}

\firstline{Bahuṁ ve saraṇaṁ yanti}

\begin{leader}
  [Ha꜓nda mayaṁ khemākhema-sa꜕raṇa-gamana-\\
  -pa꜕ridīpikā-gāthā꜓yo bha꜕ṇāmase]
\end{leader}

\begin{english}
  [Cantiamo ora i versi che distinguono i veri dai falsi rifugi]
\end{english}

\begin{twochants}
  Bahuṁ ve sa꜕ra꜓ṇaṁ yanti꜕ & pa꜕bba꜕tāni va꜕nāni꜓ ca \\
  Ārāma-rukkha꜕-cetyāni & manussā꜓ bha꜕ya꜕-tajji꜕tā \\
\end{twochants}

\begin{english}
  Molti rifugi cercano gli uomini spinti dalla paura –\\
  pendii montuosi e boschi,
  santuari nei parchi e alberi sacri.
\end{english}

\begin{twochants}
  N'etaṁ kho sa꜕ra꜓ṇaṁ khemaṁ & n'etaṁ sa꜕raṇam-u꜓tt꜕amaṁ \\
  N'etaṁ sa꜕raṇam-āgamma & sa꜕bba-dukkhā꜓ pa꜕mucca꜕ti \\
\end{twochants}

\begin{english}
  Un tale rifugio non è sicuro,
  un tale rifugio non è supremo,
  un tale rifugio non porta
  la completa liberazione dalla sofferenza.
\end{english}

\begin{twochants}
  Yo ca꜕ Buddhañca꜕ Dhammañca꜕ & sa꜓ṅghañca꜕ sa꜓ra꜕ṇaṁ ga꜕to \\
  Ca꜕ttāri a꜕riya-saccāni & sa꜕mmappaññāya꜓ pa꜕ss꜕ati \\
\end{twochants}

\begin{english}
  Chiunque vada a rifugiarsi
  nella Triplice Gemma
  vede con retto discernimento
  le Quattro Nobili Verità:
\end{english}

\begin{twochants}
  Dukkhaṁ dukkha-sa꜕muppādaṁ & dukkhassa ca꜕ a꜕ti꜕kka꜕maṁ \\
  A꜕riyañ-c'a꜕ṭṭh'a꜓ṅgi꜕kaṁ maggaṁ & dukkhūpasa꜕ma꜕-gāmi꜓naṁ \\
\end{twochants}

\begin{english}
  La sofferenza e la sua origine
  e ciò che si trova al di là –
  il Nobile Ottuplice Sentiero
  che conduce alla fine della sofferenza.
\end{english}

\begin{twochants}
  Etaṁ kho sa꜕ra꜓ṇaṁ khemaṁ & etaṁ sa꜕raṇam-u꜓tta꜕maṁ \\
  Etaṁ sa꜕raṇam-āgamma & sa꜕bba-dukkhā꜓ pa꜕mucca꜕ti \\
\end{twochants}

\begin{english}
  Un tale rifugio è sicuro,
  un tale rifugio è supremo,
  un tale rifugio porta veramente
  la completa liberazione da ogni sofferenza.
\end{english}

\chapter{Versi sulle ricchezze di un Nobile}

\firstline{Yassa saddhā tathāgate}

\begin{leader}
  [Ha꜓nda mayaṁ a꜕riya-dhana-gāthā꜓yo bha꜕ṇāmase]
\end{leader}

\begin{english}
  [Cantiamo ora i versi sulle ricchezze di un Nobile]
\end{english}

\begin{twochants}
  Yassa꜕ sa꜕ddhā tathā꜓ga꜕te & a꜕ca꜕lā su꜕pa꜕tiṭṭhi꜓tā \\
  Sī꜓lañca꜕ yassa꜕ kalyāṇaṁ & a꜕riya-kantaṁ pasa꜓ṁsi꜕taṁ \\
\end{twochants}

\begin{english}
  Colui la cui fede nel Tathāgata
  è incrollabile e ben stabilita,
  la cui virtù è bella,
  amata e lodata dai Nobili;
\end{english}

\begin{twochants}
  Sa꜓ṅghe pa꜕sā꜕do yass'atthi & uju-bhūtañca da꜓ss꜕anaṁ \\
  A꜕daliddo't꜕i taṁ āhu꜕ & a꜕moghaṁ ta꜕ssa꜕ jīvi꜓taṁ \\
\end{twochants}

\begin{english}
  La cui fiducia è nel Saṅgha,
  che vede le cose rettamente come sono,
  si dice che non invano
  e senza inganno è la sua vita.
\end{english}

\begin{twochants}
  Tasmā sa꜕ddhañca꜕ sī꜓lañca꜕ & pasādaṁ dhamma-da꜓ssa꜕naṁ \\
  A꜕nuyuñjetha medhāvī & sa꜕raṁ buddhāna sā꜓sa꜕naṁ \\
\end{twochants}

\begin{english}
  Alla virtù e alla fede,
  alla fiducia di vedere la verità,
  a queste i saggi si dedicano,
  l'insegnamento del Buddha nella loro mente.
\end{english}

\chapter{Versi sulle tre caratteristiche}

\firstline{Sabbe saṅkhārā aniccā'ti}

\begin{leader}
  [Ha꜓nda mayaṁ ti-lakkhaṇ'ādi-gāthā꜓yo bha꜕ṇāmase]
\end{leader}

\begin{english}
  [Cantiamo ora i versi sulle tre caratteristiche]
\end{english}

\begin{twochants}
  Sa꜕bbe sa꜓ṅkhā꜓rā a꜕ni꜓ccā't꜕i & yadā paññāya꜓ pa꜕ssa꜕ti \\
  Atha nibbinda꜕ti dukkhe & esa꜕ maggo vi꜓su꜕ddh꜓iyā \\
\end{twochants}

\begin{english}
  ‘Impermanenti sono tutte le cose condizionate’ –
  quando con saggezza questo si vede
  ci si stanca di ogni dukkha;
  questo è il sentiero verso la purezza.
\end{english}

\begin{twochants}
  Sa꜕bbe sa꜓ṅkhā꜓rā du꜕kkhā't꜕i & yadā paññāya꜓ pa꜕ssa꜕ti \\
  Atha nibbinda꜕ti dukkhe & esa꜕ maggo vi꜓su꜕ddh꜓iyā \\
\end{twochants}

\begin{english}
  ‘Dukkha sono tutte le cose condizionate’ –
  quando con saggezza questo si vede
  ci si stanca di ogni dukkha;
  questo è il sentiero verso la purezza.
\end{english}

\begin{twochants}
  Sa꜕bbe dhammā ana꜓ttā'ti꜕ & yadā paññāya꜓ pa꜕ssa꜕ti \\
  Atha nibbinda꜕ti dukkhe & esa꜕ maggo vi꜓su꜕ddh꜓iyā \\
\end{twochants}

\begin{english}
  ‘Non c'è un sé in alcuna cosa’ –
  quando con saggezza questo si vede
  ci si stanca di ogni dukkha;
  questo è il sentiero verso la purezza.
\end{english}

\clearpage

\begin{twochants}
  A꜕ppa꜕kā te manusse꜓su꜕ & ye janā pāra-gāmi꜓no \\
  A꜕thāyaṁ i꜕ta꜕rā pajā & tīram-evānudhā꜓va꜕ti \\
\end{twochants}

\begin{english}
  Pochi tra gli esseri umani
  sono coloro che vanno oltre,
  eppure ci sono le molte genti
  che vagano sempre su questa sponda.
\end{english}

\begin{twochants}
  Ye ca꜕ kho sammad-akkhāte & dhamme dhammānuva꜓tt꜕ino \\
  Te ja꜕nā pā꜕ram-essanti & ma꜕ccu-dheyyaṁ sudu꜓tta꜕raṁ \\
\end{twochants}

\begin{english}
  Ovunque il Dhamma sia ben insegnato,
  coloro che si addestrano in linea con esso
  sono quelli che attraverseranno
  il regno della morte così difficile da fuggire.
\end{english}

\begin{twochants}
  Kaṇhaṁ dhammaṁ vi꜕ppahā꜓ya & su꜕kkaṁ bhāvetha꜕ paṇḍi꜓to \\
  Okā a꜕noka꜕m-āgamma & viveke ya꜕tth꜕a dūramaṁ \\
  Ta꜕trābh꜕irat꜕im-iccheyya & hi꜕tvā kāme a꜕kiñc꜓ano \\
\end{twochants}

\begin{english}
  Abbandonando gli stati oscuri,
  i saggi perseguono la luce;
  dalle inondazioni raggiungono la terra asciutta
  vivendo ritirati, cosa così difficile da fare.
  Una tale rara delizia si dovrebbe desiderare,
  i piaceri dei sensi gettati via,
  non avendo nulla.
\end{english}
